% ===== 原文 pp.103--108 =====
% [原文 p.103]

\begin{itemize}[leftmargin=1.4em,itemsep=4pt,topsep=1pt]
\item \textbf{双荷真空：} 在双荷真空中，UV 线表现出以下行为：
\begin{equation}
\begin{aligned}
  \cW^{UV}_j&\longrightarrow\text{禁闭}\,, && j\notin\Z\,,\\
  \cW^{UV}_j&\longrightarrow\text{去禁闭}\,, && j\in\Z\,,\\
  \cH^{UV}_n&\longrightarrow\text{去禁闭}\,, && n\in 2\Z\,,\\
  \cH^{UV}_n&\longrightarrow\text{禁闭}\,, && n\notin 2\Z\,,\\
  \cD^{UV}_{j,n}&\longrightarrow\text{去禁闭}\,, && 2j-n\in 2\Z\,,\\
  \cD^{UV}_{j,n}&\longrightarrow\text{禁闭}\,, && 2j-n\notin 2\Z\,.
\end{aligned}
\tag{4.228}
\end{equation}
\end{itemize}

\para{$\SU(2)$ 理论。}

现在让我们为 UV 规范理论取规范群
\begin{equation}
  \cG=\SU(2)\,, \tag{4.229}
\end{equation}
并研究 IR 中会发生什么。带电的真线算符是
\begin{equation}
  \cD^{UV}_{j,n}\,,\qquad j\notin\Z,\ n\in 2\Z\,, \tag{4.230}
\end{equation}
它们在两个真空中都禁闭。于是我们得知如下结论。

\begin{statementbox}{陈述 4.11：4d $\cN=1$ $\SU(2)$ SYM 中的 1-形式 SSB}
4d $\cN=1$ $\SU(2)$ SYM 理论的 $\Z_2$ 1-形式对称性在两个真空中都自发地\textbf{未破缺}。
\end{statementbox}

两个真空中的 IR 理论都应当是一个带 $\Z_2$ 1-形式对称性的 4d TQFT。为了确定这个 TQFT，我们需要研究在 IR 中存活下来的算符：

\begin{itemize}[leftmargin=1.4em,itemsep=4pt,topsep=1pt]
\item \textbf{磁单极子真空：} 我们需要考虑线算符
\begin{equation}
  \cD^{UV}_{j,n}\,,\qquad j\in\Z,\ n\in\Z\,, \tag{4.231}
\end{equation}
以及产生 $\Z_2$ 1-形式对称性的面算符
\begin{equation}
  U_2\,, \tag{4.232}
\end{equation}
真线算符是 $n\in 2\Z$ 的那些，它们在 1-形式对称性下都不带电。这意味着 1-形式对称性在 IR 中必须是非忠实的，即 $U_2$ 在 IR 中必须同构于恒等面算符。这可由以下事实推出：TQFT 中任何非恒等的真拓扑算符都必须与该 TQFT 中某个别的非恒等真拓扑算符发生非平凡的编织。既然 $U_2$ 不与任何线发生非平凡编织，它就必定是恒等算符。这意味着 (4.231) 中所有线算符都是真线，但没有非恒等的面算符可供它们编织。因此所有 IR 线算符也都必须是恒等的。于是我们得到 IR 中一个只携带恒等算符的完全平凡的理论。

% [原文 p.104]

到目前为止，我们还没有把 $\Z_2$ 1-形式对称性的耦合纳入 IR 理论。由于 4d 中 $\Z_2$ 1-形式对称性存在非平凡的 SPT 相 (4.161)，这里有两种可能性。为了确定是哪一种，我们需要在 UV 中为 $\Z_2$ 1-形式对称性选择一个耦合，即选择扭曲扇区中非平凡带荷的线是
\begin{equation}
  \cD^{UV}_{j,n}\,,\qquad j\notin\Z,\ n\notin 2\Z\,, \tag{4.233}
\end{equation}
还是
\begin{equation}
  \cD^{UV}_{j,n}\,,\qquad j\in\Z,\ n\notin 2\Z\,. \tag{4.234}
\end{equation}
按惯例，选择 (4.233) 就直接称为 $\SU(2)$ 理论，而选择 (4.234) 称为 $T\SU(2)$ 理论。对 $\SU(2)$ 理论，没有任何带荷的线在 IR 中存活，因此我们得到平凡的 SPT 相。另一方面，对 $T\SU(2)$ 理论，我们确实在 IR 的扭曲扇区中有带荷的线，因此得到非平凡的 SPT 相。于是带 $\Z_2$ 1-形式对称性的 IR 4d TQFT 为
\begin{equation}
\begin{aligned}
  \cT^M_{SU(2)}&=\text{平凡}\,,\\
  \cT^M_{TSU(2)}&=\text{SPT}\,,
\end{aligned}
\tag{4.235}
\end{equation}
其中 SPT 表示非平凡的 SPT 相，上标 $M$ 表示磁单极子真空。
\item \textbf{双荷真空：} 需要纳入考虑的线算符是
\begin{equation}
  \cD^{UV}_{j,n}\longrightarrow\text{去禁闭}\,,\qquad 2j-n\in 2\Z\,. \tag{4.236}
\end{equation}
真线算符满足 $j\in\Z$、$n\in 2\Z$，它们在 1-形式对称性下不带电。因此 $U_2$ 必须又是恒等算符，并且与磁单极子真空的情形类似，所有线算符也都必须是恒等的。因此，若不把 $\Z_2$ 1-形式对称性考虑进来，IR 理论是平凡的。

在把 $\Z_2$ 1-形式对称性考虑进来之后，4d TQFT 为
\begin{equation}
\begin{aligned}
  \cT^D_{SU(2)}&=\text{SPT}\,,\\
  \cT^D_{TSU(2)}&=\text{平凡}\,,
\end{aligned}
\tag{4.237}
\end{equation}
其中上标 $D$ 表示双荷真空。
\end{itemize}

% [Tutor Note: p.103--104 (4.235)、(4.237) 的判据是：IR 中是否存在带荷的扭曲扇区线决定 SPT 相（存在 → 非平凡 SPT）。以磁单极子真空为例：$\SU(2)$ 理论取耦合 (4.233)，其带电扭曲扇区线为 $j\notin\Z,n\notin 2\Z$，由 (4.227) 第 4 行（$j\notin\Z,n\in\Z$ 时禁闭）全部禁闭，故 $\cT^M_{\SU(2)}=$ 平凡；而 $T\SU(2)$ 理论取 (4.234)，其带电扭曲扇区线 $j\in\Z,n\notin 2\Z$ 由 (4.227) 第 5 行去禁闭、在 IR 中存活，故 $\cT^M_{T\SU(2)}=$ SPT。双荷真空同理可用 (4.228) 的第 5、6 行核对（$n\notin 2\Z$ 即 $n$ 为奇数，$2j-n$ 的奇偶性决定去禁闭/禁闭）。$M$、$D$ 只标记真空，不是对称性指标。]

% [原文 p.105]

传统上，携带平凡 SPT 的真空称为\textbf{禁闭真空}，而非平凡的 SPT 则称为\textbf{斜禁闭真空}。

注意 $\SO(3)_+$ 与 $T\SO(3)_+$ 理论的定义是这样选取的：对 $\SU(2)$ 理论的 $\Z_2$ 1-形式对称性作规范化给出的是 $\SO(3)_+$ 理论（而不是 $T\SO(3)_+$ 理论）。

\para{$\SO(3)_+$ 理论。}

现在取规范群
\begin{equation}
  \cG=\SO(3)\,, \tag{4.238}
\end{equation}
且离散 theta 角平凡。带电的真线算符是
\begin{equation}
  \cD^{UV}_{j,n}\,,\qquad j\in\Z,\ n\notin 2\Z\,, \tag{4.239}
\end{equation}
它们在磁单极子真空中去禁闭，但在双荷真空中禁闭。于是我们得知如下结论。

\begin{statementbox}{陈述 4.12：4d $\cN=1$ $\SO(3)_+$ SYM 中的 1-形式 SSB}
4d $\cN=1$ $\SO(3)_+$ SYM 理论的 $\Z_2$ 1-形式对称性在一个真空中自发\textbf{破缺}，在另一个真空中自发\textbf{未破缺}。
\end{statementbox}

让我们研究两个真空中的 IR TQFT。

\begin{itemize}[leftmargin=1.4em,itemsep=4pt,topsep=1pt]
\item \textbf{磁单极子真空：} 由于我们有去禁闭的带荷真线 (4.239)，$U_2$ 是一个非恒等的面算符。类似地，(4.239) 中所有线都是非恒等的拓扑算符。然而这些线在 IR 中全都是同一条线。这是因为，若我们有两条不同的线 $L_1$ 与 $L_2$，则 $L_1^{-1}L_2$ 将是一条非恒等的线，它不与任何面发生非平凡编织，从而导致矛盾。我们把这条 IR 线算符记为 $U_1$。于是，作为去禁闭的结果，我们在 IR 中得到了一个\textbf{涌现的 2-形式对称性}！我们可以认出 $U_2$ 与 $U_1$ 的性质正是 4d $\Z_2$ 规范理论
\begin{equation}
  \pi\int a_1\cup\delta b_2 \tag{4.240}
\end{equation}
中面算符与线算符的性质；这个 4d $\Z_2$ 规范理论就是 IR 有效理论
\begin{equation}
  \cT^M_{\SO(3)_+}=\text{$\Z_2$ 规范理论}\,. \tag{4.241}
\end{equation}
我们可以修改 UV 中的 1-形式耦合、从 $T\SO(3)_+$ 理论出发，但这只影响扭曲扇区线，而它们全都禁闭，因此该耦合不影响 IR 物理
\begin{equation}
  \cT^M_{T\SO(3)_+}=\text{$\Z_2$ 规范理论}\,. \tag{4.242}
\end{equation}
\end{itemize}

% [Tutor Note: p.105 (4.240) 中 $a_1$ 与 $b_2$ 分别是 $\Z_2$ 1-形式与 2-形式规范场，$\pi\int a_1\cup\delta b_2$ 是 $\Z_2$ 规范理论的拓扑作用量。核对方法：$a_1,b_2$ 取值于 $\Z_2=\{0,1\}$，故相位为 $\exp(\ii\pi\int a_1\cup\delta b_2)=(-1)^{\int a_1\cup\delta b_2}\in\{\pm1\}$，与 $U_2$、$U_1$ 的编织相位相符——这正是 $\Z_2$ 规范理论中 Wilson 线与 't Hooft 线的非平凡编织。]

% [原文 p.106]

\begin{itemize}[leftmargin=1.4em,itemsep=4pt,topsep=1pt]
\item \textbf{双荷真空：} 这里的情况与上面讨论的 $\SU(2)$ 各情形类似，人们发现
\begin{equation}
\begin{aligned}
  \cT^D_{\SO(3)_+}&=\text{SPT}\,,\\
  \cT^D_{T\SO(3)_+}&=\text{平凡}\,,
\end{aligned}
\tag{4.243}
\end{equation}
其中 $\SO(3)_+$ 的 1-形式耦合这样选取，使带荷的扭曲扇区线为
\begin{equation}
  \cD_{j,n}\,,\qquad j\notin\Z,\ n\notin 2\Z\,, \tag{4.244}
\end{equation}
而对 $T\SO(3)_+$，带荷的扭曲扇区线为
\begin{equation}
  \cD_{j,n}\,,\qquad j\notin\Z,\ n\in 2\Z\,. \tag{4.245}
\end{equation}
\end{itemize}

\para{$\SO(3)_-$ 理论。}

分析与 $\SO(3)_+$ 的情形类似，留给感兴趣的读者。最终结果为
\begin{equation}
\begin{aligned}
  \cT^M_{\SO(3)_-}&=\text{SPT}\,,\\
  \cT^M_{T\SO(3)_-}&=\text{平凡}\,,\\
  \cT^D_{\SO(3)_-}&=\text{$\Z_2$ 规范理论}\,,\\
  \cT^D_{T\SO(3)_-}&=\text{$\Z_2$ 规范理论}\,.
\end{aligned}
\tag{4.246}
\end{equation}

注意 $\SO(3)_-$ 与 $T\SO(3)_-$ 理论的定义是这样选取的：对 $T\SU(2)$ 理论的 $\Z_2$ 1-形式对称性作规范化给出的是 $\SO(3)_-$ 理论（而不是 $T\SO(3)_-$ 理论）。

\begin{notebox}{禁闭的非拉格朗日理论}
在上面关于 4d $\cN=1$ SYM 中禁闭/去禁闭的研究中，拉格朗日描述对于得到所呈现的结果是关键的。然而，禁闭与去禁闭的定义并不需要任何拉格朗日描述。这就提出了一个问题：如何在非拉格朗日或强耦合的 4d 理论中研究禁闭/去禁闭？最近 [68] 对这样一类理论给出了该问题的一个答案：这些理论允许用 M-理论中的 M5 膜、或作为 6d $\cN=(2,0)$ SCFT 的紧致化来构造，其中也给出了禁闭非拉格朗日理论的例子。
\end{notebox}

\para{$S$ 与 $T$ 操作。}

上面考虑的所有 UV 理论都通过规范化相互联系。我们已经讨论过两个这样的规范化
\begin{equation}
\begin{aligned}
  \SU(2)/\Z_2^{(1)}&=\SO(3)_+\,,\\
  T\SU(2)/\Z_2^{(1)}&=\SO(3)_-\,.
\end{aligned}
\tag{4.247}
\end{equation}

% [原文 p.107]

习惯上把规范化记为 $S$。用它来表示，上面的规范化写成关系式
\begin{equation}
\begin{aligned}
  S(\SU(2))&=\SO(3)_+\,,\\
  ST(\SU(2))&=\SO(3)_-\,.
\end{aligned}
\tag{4.248}
\end{equation}
因此，$S$ 是规范化的操作，而 $T$ 是叠加一个 SPT 相的操作。可以证明 [4,93]，这些操作满足
\begin{equation}
  (ST)^3=1\,, \tag{4.249}
\end{equation}
这也是生成 $SL(2,\Z)$ 群的 $S$ 与 $T$ 矩阵所满足的恒等式。出于这个原因，上述 $S$ 与 $T$ 操作被称为“对带 $\Z_2$ 1-形式对称性的 4d QFT 的 $SL(2,\Z)$ 操作”。应当注意，在我们这个语境中 $T$ 是二阶操作，而不是通常的无穷阶操作。把它与“规范化两次等同于解除规范化”这一事实，即
\begin{equation}
  S^2=1\,, \tag{4.250}
\end{equation}
结合起来，我们发现这些 $S$ 与 $T$ 变换构成的群是 $PSL(2,\Z_2)$。

利用关系 (4.249) 连同 (4.247)，可以推出完整的规范化之网。例如，若想知道对 $T\SO(3)_-$ 理论的 $\Z_2$ 1-形式对称性作规范化的结果，就是在计算
\begin{equation}
  S(T\SO(3)_-)=STST(\SU(2))=TS(\SU(2))=T\SO(3)_+\,. \tag{4.251}
\end{equation}
或者，也可以用 1-形式规范化下线算符带电扇区与扭曲扇区的交换，来推出规范化后的理论。

\para{UV 与 IR 规范化的相容性。}

人们也可以考虑对上面讨论的 IR TQFT 的 1-形式对称性作规范化。我们有
\begin{equation}
\begin{aligned}
  S(\text{平凡})&=\text{$\Z_2$ 规范理论}\,,\\
  S(\text{SPT})&=\text{SPT}\,,\\
  S(\text{$\Z_2$ 规范理论})&=\text{平凡}\,.
\end{aligned}
\tag{4.252}
\end{equation}
这可以通过交换线算符的扭曲扇区与带电扇区而容易地看出：

\begin{itemize}[leftmargin=1.4em,itemsep=4pt,topsep=1pt]
\item 在平凡理论中，我们有一条不带电的扭曲扇区线。它不过是位于 $U_2$ 端点处的恒等线，而如我们所讨论的，$U_2$ 不过是恒等面算符。规范化之后，这条线给出未扭曲扇区中的一条带荷线，这意味着规范化后的理论是 $\Z_2$ 规范理论。
\end{itemize}

% [Tutor Note: p.107 (4.251) 的第二步 $STST=TS$ 可由 (4.249)--(4.250) 推出：$(ST)^3=1\Rightarrow STST=(ST)^{-1}=T^{-1}S^{-1}=TS$（用到 $S^2=T^2=1$）。这正是原文所谓用 (4.249) 与 (4.247) 推出整个规范化之网的含义。可核对：$T$ 的二阶性即 $T^2=1$ 来自“叠加两次 SPT 相 = 平凡”。]

% [原文 p.108]

\begin{itemize}[leftmargin=1.4em,itemsep=4pt,topsep=1pt]
\item 在 SPT 理论中，我们有一条带电的扭曲扇区线。规范化之后，它仍留在扭曲扇区中，并且仍然带电。因此，规范化后的理论仍然是 SPT。
\item 在 $\Z_2$ 规范理论中，我们有一条带电的未扭曲扇区线。规范化之后，它变成扭曲扇区中一条不带电的线。因此，规范化后的理论是平凡理论。
\end{itemize}

$T$ 操作为
\begin{equation}
\begin{aligned}
  T(\text{平凡})&=\text{SPT}\,,\\
  T(\text{SPT})&=\text{平凡}\,,\\
  T(\text{$\Z_2$ 规范理论})&=\text{$\Z_2$ 规范理论}\,.
\end{aligned}
\tag{4.253}
\end{equation}
现在，$S$ 与 $T$ 变换应当与 RG 流对易。确实，这可以显式地检验。例如，我们应当有
\begin{equation}
  S\left(\cT^M_{SU(2)}\right)=\cT^M_{S(SU(2))}=\cT^M_{\SO(3)_+}\,, \tag{4.254}
\end{equation}
而事实正是如此。作为另一个例子，我们应当有
\begin{equation}
  S\left(\cT^D_{T\SO(3)_-}\right)=\cT^D_{T\SO(3)_+}\,, \tag{4.255}
\end{equation}
读者可以容易地看出它也是对的。

\subsection{对称性拓扑场论（Symmetry TFTs）}

我们已经看到，离散对称性的任何规范化都可以解除，只需随后对对偶对称性作一次规范化即可。这意味着离散规范化并不改变一个理论中可观测量的全部信息。然而，正如我们已经看到的，离散规范化确实会把不同类型的可观测量彼此混合，例如 4.3.4 节讨论的扭曲扇区与带电扇区的交换。其结果是，同一个物理可观测量在与某个离散对称性相关联的不同\emph{规范框架}中可能以不同的面貌出现。对称性拓扑场论（TFT），缩写为 \textbf{SymTFT} 或 \textbf{Symmetry TFT}，是把物理可观测量与其在不同规范框架下的行为解耦的有力工具。它们还编码该对称性的 't Hooft 反常以及可能的高阶群结构。关于 SymTFT 的一般逻辑，我们请读者参阅 [179]；关于可逆对称性语境下 SymTFT 的详细讨论，请参阅 [4,102,180]。最近，[60] 中出现了关于不可逆对称性的 SymTFT 的详细讨论，其中也阐明了 [179] 的一般概念。

% [Tutor Note: p.108 是本块内 §4.5 的起始，只给出 SymTFT 的动机：离散规范化之间的“规范化/解除规范化”（$S$ 操作与其逆）意味着同一物理可观测量的不同规范框架描述，SymTFT 就是把这种框架依赖性解耦的拓扑场论；它同时编码 't Hooft 反常与高阶群结构。原文在此并列给出两种写法：symmetry topological field theory（TFT）与 SymTFT/Symmetry TFT，本译本统一作“对称性拓扑场论（SymTFT）”。原文给出的文献：[179]（一般逻辑）、[4,102,180]（可逆对称性的具体讨论）、[60]（不可逆对称性）。]
